\section{Graph-construction definitions}
\label{sec:s-constructions}

The IBM construction study uses the edge-aware SAGE-derived h64 classifier as its reference. Every intervention retains variant, protocol, seed, and optimizer schedule unless the intervention itself changes the training structure. Table~\ref{tab:s-construction-cards} records the operation and question for each matched candidate.

\begin{table}[tbp]
\centering
\footnotesize
\setlength{\tabcolsep}{3.2pt}
\renewcommand{\arraystretch}{1.12}
\caption{Matched IBM graph-construction and feature interventions.}
\label{tab:s-construction-cards}
\begin{tabularx}{\textwidth}{@{}>{\raggedright\arraybackslash}p{0.16\textwidth} Y Y >{\raggedright\arraybackslash}p{0.20\textwidth}@{}}
\toprule
\textbf{Construction} & \textbf{Operation} & \textbf{Question} & \textbf{Matched scope} \\
\midrule
NoEdge & zero transaction attributes; structure retained & tests edge-attribute contribution & matched to h64 reference \\
ShuffledEdge & edge attributes permuted within temporal masks & destroys edge-feature alignment while preserving marginal values & matched to h64 reference \\
DegreeOnly & endpoint-degree attributes replace transaction features & tests structural summaries without original attributes & matched to h64 reference \\
DegreeCap & q=.995 cap on training structure & tests hub sensitivity and resource/performance tradeoff & matched to h64 reference \\
RecentWindow & most recent 50\% of training edges & tests recency and reduced computation & matched to h64 reference \\
Sender-receiver alias & Restates the already materialized sender-to-receiver transaction edges & Contract identity check; no new transformation & Diagnostic only; never counted as an independent method \\
\bottomrule
\end{tabularx}
\vspace{2pt}
\begin{minipage}{\textwidth}\footnotesize\emph{Scope and reading.} Every intervention retains variant, protocol, seed, optimizer schedule, and h64 reference family unless the operation itself changes training structure. The sender-receiver row is an alias, not replication.\end{minipage}
\end{table}


\paragraph{Feature controls.}
NoEdge sets all eight transaction attributes to zero while preserving account structure and histories. ShuffledEdge permutes transaction-feature rows separately within train, validation, and test masks, preserving marginal distributions and split sizes while breaking transaction-to-feature alignment. DegreeOnly replaces the transaction attributes with source/destination in/out-degree summaries computed from training edges.

\paragraph{Structural controls.}
DegreeCap removes training edges incident to accounts above the 0.995 active-training-degree quantile, subject to safety guards. RecentWindow uses only the chronologically latest 50\% of training edges to construct the training structure, changing both represented history and computation.

\paragraph{Identity diagnostic.}
The sender-receiver row records that the materialized arrays already encode sender-to-receiver transactions. It performs no independent graph transformation. Its matched zero rows remain useful as an identity check but cannot count as replication or as another method in a winner set.

\paragraph{Interpretive limit.}
These interventions isolate implementation-level representation and construction choices within the recorded IBM contracts. They do not identify causal effects for arbitrary AML networks, and Medium GINE remains outside the feasible set.
